% =============================================================
% layout.tex — B5 紙面レイアウト / ヘッダ・フッタ / 見出し / 目次
% =============================================================

% ── 日本語 ──
\usepackage{luatexja}
\usepackage{luatexja-fontspec}

% ── ページ余白 ──
\usepackage[b5paper,
  top=25mm, bottom=25mm,
  inner=22mm, outer=18mm,
  headheight=14pt, headsep=10mm, footskip=14mm
]{geometry}

% ── パッケージ群 ──
\usepackage{fancyhdr}
\usepackage{enumitem}
\usepackage{parskip}
\usepackage{graphicx}
\usepackage{adjustbox}
\usepackage{array}
\usepackage{booktabs}
\usepackage{longtable}
\usepackage{tabularx}
\usepackage{colortbl}
\usepackage{multicol}
\usepackage{tikz}
\usetikzlibrary{positioning,arrows.meta,shapes.geometric,fit,backgrounds,calc,shadows.blur}
\usepackage[most]{tcolorbox}
\tcbuselibrary{skins,breakable,xparse}
\usepackage{hyperref}
\usepackage{amssymb}
\usepackage{titletoc}
\usepackage{etoolbox}
\usepackage{xparse}
\usepackage{pgffor}

% ── ハイパーリンク ──
\hypersetup{
  colorlinks=true,
  linkcolor=Navy,
  urlcolor=Teal,
  bookmarksnumbered=true,
  pdftitle={図でつながる ITパスポート},
}

% =============================================================
% 章扉デザイン: jlreq では \chapter をそのまま使い、
% 装飾は \chaptermeta 等のコマンドで付加する方針とする。
% =============================================================

% =============================================================
% section / subsection デザイン（jlreq 互換）
% =============================================================

% =============================================================
% 目次デザイン
% =============================================================
\setcounter{tocdepth}{1}
\renewcommand{\contentsname}{目次}

\titlecontents{chapter}[0pt]
  {\addvspace{10pt}}
  {\large\bfseries\color{Navy}\makebox[2.2em][l]{\thecontentslabel}}
  {\large\bfseries\color{Navy}}
  {\normalfont\color{Muted}\ \titlerule*[0.7pc]{\footnotesize.}\ \contentspage}
  [\vspace{1pt}{\color{HairLine}\rule{\linewidth}{0.3pt}}]

\titlecontents{section}[2em]
  {\addvspace{2pt}\small}
  {\color{SubBlue}\makebox[2em][l]{\thecontentslabel}}
  {}
  {\normalfont\footnotesize\color{Muted}\ \titlerule*[0.7pc]{\tiny.}\ \contentspage}

% =============================================================
% ヘッダ / フッタ
% =============================================================
\pagestyle{fancy}
\fancyhf{}
\fancyhead[LE]{\small\color{Muted}\sffamily 図でつながる ITパスポート}
\fancyhead[RO]{\small\color{SubBlue}\nouppercase\leftmark}
\fancyfoot[C]{\small\color{Muted}\thepage}
\renewcommand{\headrulewidth}{0pt}
\renewcommand{\footrulewidth}{0pt}
\renewcommand{\headrule}{%
  \hbox to\headwidth{\color{HairLine}\leaders\hrule height 0.3pt\hfill}%
}
\makeatletter
\def\chaptermark#1{\markboth{第\,\thechapter\,章\quad #1}{}}
\makeatother

% ── リスト / 段落 ──
\setlist[itemize]{leftmargin=1.4em, itemsep=2pt, topsep=4pt}
\setlist[enumerate]{leftmargin=1.6em, itemsep=2pt, topsep=4pt}
\setlength{\parindent}{1em}
\setlength{\parskip}{0.4em plus 0.1em minus 0.05em}
\renewcommand{\baselinestretch}{1.08}
